% RESEARCH_DRAFT_SHA256: 73fc723d91d293af4a4b2b0fee5a4478c8147f0409b3dc46800463b952f8783c
% Layout and source-reference wrappers expanded for the web renderer.
\documentclass[11pt]{article}
\usepackage[margin=25mm]{geometry}
\usepackage{amsthm}
\usepackage{fontspec,unicode-math}
\setmainfont{Latin Modern Roman}
\setmathfont{Latin Modern Math}
\usepackage{amsmath,amssymb,mathtools}
\usepackage{xurl,hyperref}
\hypersetup{colorlinks=true,urlcolor=blue,linkcolor=blue}
\setcounter{secnumdepth}{0}
\setlength{\parindent}{0pt}
\setlength{\parskip}{5pt}
\setlength{\emergencystretch}{3em}
\newcommand{\sref}[2]{\hyperlink{source:#1:#2}{[S#2]}}
\newcommand{\cataloguescope}[1]{\paragraph{Source question.}#1}
\newcommand{\floor}[1]{\left\lfloor#1\right\rfloor}
\newtheorem{proposition}{Proposition}
\newtheorem{lemma}[proposition]{Lemma}
\begin{document}
\paragraph{Research manuscript.} OpenAI \texttt{gpt-6-astra}, reasoning \texttt{ultra}; 11 October 2026.
\paragraph{Preserved source body.} The complete catalogue mathematical text and sources follow with layout and source links expanded, including source qualifications and its historical preparation notes. New research begins at the section headed Research attempt.

% UnsolvedMath ID 30006104; source catalogue code OWR-14298804-004
\section{Integer Feasibility with Bounded Subdeterminants}\label{30006104}
{\small UnsolvedMath ID 30006104 · Computer Science}
\subsection{Definitions and mathematical statement}

\paragraph{Shared notation.}Standard number systems are denoted by $\mathbb N,\mathbb Z,\mathbb Q,\mathbb R,\mathbb C$; index ranges, computational models and specialized notation follow the cited problem source and the definitions below.


A full-rank subdeterminant is the determinant of a square submatrix of order equal to the rank of $A$. Integer feasibility asks whether the constraint system has a vector in $\mathbb Z^n$.

Polynomial time means a number of elementary steps bounded by a polynomial in the encoded input length. A decision problem is in $\mathrm{NP}$ if yes-instances have polynomial-length certificates verifiable in polynomial time; NP-hardness is under polynomial-time reductions.
\paragraph{Statement recorded in the supplied source.}
Let $P=\{x\in\mathbb R^n:Ax\leq b\}$, where $A$ is integral and the absolute values of its full-rank subdeterminants are bounded by $\Delta$. For arbitrary $\Delta\geq3$, can one decide in polynomial time whether $P$ contains an integer vector? \href{https://doi.org/10.4171/owr/2024/50}{[S2]}
\paragraph{Scope qualification.}
The question does not specify whether $\Delta$ is fixed or encoded as part of the input. Full-rank determinant bounds are not replaced by bounds on every smaller square minor.
\paragraph{Source status.}
The supplied catalogue labels this record open. Its source-status metadata was checked against the supplied snapshot on 2026-10-11. This import does not independently establish current open status or mathematical correctness. The catalogue literature assessment is dated 2026-08-22.
\subsection{Short English statement}
Decide integer feasibility efficiently when full-rank determinants are bounded by the stated parameter.
\subsection{Sources}
\begin{itemize}
\item[S1] Ulam AI and the UnsolvedMath Contributors, UnsolvedMath: A Curated Collection of Open Mathematics Problems, record 30006104 (OWR-14298804-004). Catalogue curation and metadata: CC BY 4.0. The supplied source question is retained; definitions and the English summary are editorial additions.
\url{https://huggingface.co/datasets/ulamai/UnsolvedMath}.
\item[S2] Combinatorial Optimization (2025), Oberwolfach Reports, DOI 10.4171/owr/2024/50, pp. 2949--3000
\url{https://doi.org/10.4171/owr/2024/50}.
\item[S3] Kuhlmann et al. in Oberwolfach Report 50
\url{https://doi.org/10.4171/OWR/2024/50}.
\item[S4] Celaya et al., Bounded non-vanishing subdeterminants
\url{https://doi.org/10.1016/j.jcta.2024.106003}.
\textit{Underlying Oberwolfach Reports material retains its source-specific licence, including CC BY-SA 4.0 where applicable; see THIRD\_PARTY\_NOTICES.md.}
\end{itemize}
\paragraph{Preparation and review.}
Original question attributed to the cited source and UnsolvedMath contributors. Codex (GPT-6; exact serving revision not exposed) assisted with definitions, English summaries, source inspection and formatting. No new human mathematical authorship or independent proof verification is asserted. Canonical ID 30006104 and source code OWR-14298804-004 checked against the complete supplied registry and existing repository statements; known flagged repetitions and obvious identical conjectures excluded. Full mathematical-equivalence/literature review has not been performed. Statement preparation and final review followed contribution prompts C and D; automated checks verify structure and identity, not open status.
\label{end:30006104}

\section{Research attempt}
\subsection{Literature, source context and scope}
The full catalogue body, including qualifications and historical source status, is preserved above. The original workshop report was inspected in the locally cached publisher PDF text. This attempt distinguishes the source question from any explicitly restricted model introduced below. Repeated records retain their own source IDs and do not become independent conjectures. Elementary results and attributed literature results are not claimed as new original-target progress.

\subsection{Source convention and attempted proximity route}
The primary report explicitly calls the determinant bound a constant and assumes both $A$ and $b$ integral. We analyze that fixed-$\Delta$ setting without changing the more qualified catalogue wording. A possible enumeration route is to find an LP vertex and search an axis-aligned neighborhood whose radius depends only on $\Delta$ and the dimension. The following full-dimensional example defeats that unconditioned proximity claim.
For an integer $M\ge1$, put
\[
 U_M=\begin{pmatrix}1&M\\0&3\end{pmatrix},\qquad
 P_M=\{x\in\mathbb R^2:0\le U_Mx\le(1,1)^T\}.
\]
The four-row inequality matrix is formed by the two rows of $U_M$ and their negatives.
\begin{proposition}
Every full-rank subdeterminant has absolute value zero or three. The vertex $v_M=(-M/3,1/3)$ is at $\ell_\infty$ distance at least $M/3$ from every integer point in $P_M$.
\end{proposition}
\begin{proof}
A pair of opposite rows is dependent. Any independent pair is, up to signs and ordering, the two rows of $U_M$, whose determinant is three. Since $U_M$ is invertible, $P_M$ is the inverse image of a square and has nonempty interior; $v_M$ is the inverse image of $(0,1)$. For integer $x$, the constraint $0\le3x_2\le1$ forces $x_2=0$, after which $0\le x_1\le1$ gives exactly $(0,0)$ and $(1,0)$. Their first-coordinate distances from $v_M$ are respectively $M/3$ and $1+M/3$.
\end{proof}
Thus no radius depending solely on $(n,\Delta)=(2,3)$ works around every LP vertex in the original coordinates. The large one-by-one minor $M$ also demonstrates why a bound on full-rank minors cannot be silently replaced by a bound on every minor. Integer-preserving coordinate preprocessing could evade this obstruction, and this family is easy to solve directly.
\subsection{Verification and remaining gaps}
For $M=30$, all six two-row determinants were computed and the two integer feasible points checked; the displayed vertex has distance ten from the nearest one. A general algorithm may choose different vertices or transformed coordinates. Neither feasibility hardness nor a new bounded-determinant algorithm follows.

\section{References and source attribution}
The primary Oberwolfach workshop report, \url{https://doi.org/10.4171/owr/2024/50}, was inspected on 11 October 2026. Relevant source passages: 2024/50 pp. 2978--2979, Kuhlmann contribution. Catalogue and contributor attribution remain in the original Sources section. Mathematical arguments labeled as this attempt are proved above; cited external results remain attributed to their authors and are not independently reproved or claimed as our discoveries.


\paragraph{Integration review.} Codex (GPT-6; exact serving revision not exposed) reviewed the written arguments and expanded layout and source-reference wrappers. Repository metadata and the production web renderer were checked using Python 3.13 and JavaScript with MathJax 3.2.2. This model-assisted review is not independent human proof verification; no PDF compilation or formal proof checking is asserted.

\section{Completion Estimate}
\noindent\textbf{COMPLETION: 0\%}
No substantiated new progress toward the entire original target is claimed. The worked results and failed route are delimited above; a known result, illustrative model, or auxiliary obstruction is not counted as original-target completion.

\section{Final status and attribution}
\textbf{Attempt status: unresolved.} This describes this research attempt, not a new ruling on the source catalogue's status or on external literature claims.
The mathematics and this self-audit were generated by OpenAI \texttt{gpt-6-astra}, reasoning \texttt{ultra}, on 11 October 2026. Codex (GPT-6; exact serving revision not exposed) orchestrated the batch and contributed formatting/repository review. No human mathematical authorship, independent human proof checking, or formal verification is claimed. The historical preparation paragraphs in the reproduced catalogue body concern that body only.

Contribution prompt A was applied to the complete definition and supplied source record, followed by the adversarial checks of prompt B. Local reproduction files are \texttt{batch011/checks.py}, \texttt{computations.json}, and the per-ID audit JSON. They were not public links at preparation time. Numerical checks are finite and do not prove unrestricted results. No LaTeX engine was found locally; compilation was not performed by this research agent.

\end{document}
